% ============================================================
%  chapters/appendix-a.tex
%  Spherepop Formalism — formal definitions and theorems
% ============================================================

\section*{Overview}

This appendix collects the formal definitions underlying the Spherepop
calculus developed in Chapters~5–9. References to main-text theorems
are given where proofs are deferred.

% ── A.1 Syntax ────────────────────────────────────────────────
\section{Syntax}

\begin{definition}{Spherepop Expression}{sp-expr}
  The set $\mathcal{E}$ of \emph{Spherepop expressions} is defined
  inductively:
  \begin{enumerate}
    \item Every \emph{atom} $a \in \mathcal{A}$ is an expression.
    \item If $e_1, \ldots, e_n \in \mathcal{E}$, then
          $\sphere{e_1, \ldots, e_n}$ is an expression
          (a \emph{sphere} containing $e_1, \ldots, e_n$).
    \item Every expression has a unique parse tree.
  \end{enumerate}
\end{definition}

\begin{definition}{Nesting Depth}{nesting-depth}
  The \emph{depth} $d : \mathcal{E} \to \mathbb{N}$ is defined by:
  \[
    d(a) = 0 \quad \text{for atoms},\qquad
    d\bigl(\sphere{e_1,\ldots,e_n}\bigr) = 1 + \max_i d(e_i)
  \]
\end{definition}

% ── A.2 Evaluation ───────────────────────────────────────────
\section{Evaluation Dynamics}

\begin{definition}{Pop Operator}{pop-op}
  The \emph{pop operator} $\pop : \mathcal{E} \to \mathcal{E} \cup \{\refuse\}$
  is defined by:
  \[
    \sphere{e_1, \ldots, e_n} \pop = \begin{cases}
      e_1 \cdots e_n & \text{if the sphere is \emph{collapsible}} \\
      \refuse        & \text{if the sphere \emph{refuses} (is protected)} \\
    \end{cases}
  \]
  where collapsibility is determined by the active constraint set $\constraint$.
\end{definition}

\begin{theorem}{Church–Rosser for Spherepop}{church-rosser}
  The evaluation relation $\to_{\mathrm{sp}}$ on $\mathcal{E}$ is
  confluent: if $e \to^* e_1$ and $e \to^* e_2$, then there exists
  $e'$ with $e_1 \to^* e'$ and $e_2 \to^* e'$.
\end{theorem}

\begin{proof}
  By induction on the structure of $e$, using the fact that the pop
  operator is deterministic on each sphere given a fixed $\constraint$.
  The key step is showing that parallel reductions commute, which follows
  from the spatial independence of non-overlapping spheres. \hfill$\square$
\end{proof}
